\documentclass{article}
\usepackage{amsmath,amssymb,amsthm,mathtools}
\newtheorem{theorem}{Theorem}
\newtheorem{lemma}{Lemma}
\newtheorem{proposition}{Proposition}
\newtheorem{corollary}{Corollary}
\newtheorem{definition}{Definition}
\newtheorem{example}{Example}
\title{ShadowBench geometry/L2/geo\_gen\_L2\_007}
\begin{document}
\maketitle
% Problem id: geometry/L2/geo_gen_L2_007
% Companion instructions: docs/instructions.md
% Candidate skeletons: docs/skeletons/
% Lean target: ShadowBench/Source/Main.lean

\begin{theorem}[brahmagupta_formula]\label{thm:brahmagupta_formula}
In Euclidean geometry, Brahmagupta's formula gives the area $K$ of a convex cyclic quadrilateral (a quadrilateral inscribed in a circle) given the lengths of its sides. 
Formally, let $A, B, C, D$ be the vertices of a convex cyclic quadrilateral in order. Let $a = |AB|, b = |BC|, c = |CD|$, and $d = |DA|$ be the lengths of the sides, and let $s$ be the semiperimeter, defined as:
\[
s = \frac{a + b + c + d}{2}
\]
Then the area $K$ of the quadrilateral is given by:
\[
K = \sqrt{(s - a)(s - b)(s - c)(s - d)}
\]
\end{theorem}
\end{document}
